\chapter{Ranks of elliptic curves}\label{ch:oc-ranks}

Volume~13 proves that $E(\Q)$ is finitely generated. It does not say how large the rank is. A random model does, and the model is not a theorem. This chapter proves the linear algebra of the model and rewrites the descent Volume~13 gave for $y^2=x^3-x$ as a Selmer group. Goldfeld's conjecture, Birch--Swinnerton-Dyer, the theorems of Bhargava--Shankar, and modularity stay stated. The orbit method is not attempted.

\section{Random maximal isotropics}

Let $p$ be a prime and let $A$ be a uniform random alternating $n\times n$ matrix over $\mathbb F_p$: zero diagonal, $A^{\mathsf T}=-A$, and the entries above the diagonal independent and uniform. The Weil pairing on a $p$-Selmer group is alternating, and the Poonen--Rains model identifies a Selmer group, in the limit of large dimension, with the kernel of such a matrix.

\begin{theorem}\label{thm:lagrangian}
\[
\mathbb E\,|\ker A|=p+1-p^{1-n}.
\]
In particular the expectation tends to $p+1$ as $n\to\infty$. The same limit is the expected size of the intersection of two random maximal isotropics in a $2n$-dimensional symplectic space, in the chart where the second is the graph of an alternating map.
\end{theorem}
\begin{proof}
The expectation $\mathbb E|\ker A|$ equals $\sum_{v\in\mathbb F_p^n}P(Av=0)$. The zero vector contributes $1$.

The symplectic group acts transitively on the nonzero vectors, and it preserves the uniform measure on alternating matrices by $A\mapsto g^{\mathsf T}Ag$. Thus $P(Av=0)$ is the same for every $v\ne0$. For $v=e_1$ the first column of $A$ below the diagonal consists of $n-1$ independent uniform entries, the diagonal entry is $0$, and the entries above the diagonal are determined by skew-symmetry. So $Ae_1$ is uniform in the hyperplane $e_1^{\perp}$, which has $p^{n-1}$ points, and
\[
P(Ae_1=0)=p^{1-n}.
\]
There are $p^n-1$ nonzero vectors, and
\[
\mathbb E\,|\ker A|=1+(p^n-1)p^{1-n}=p+1-p^{1-n}.
\]

For the geometric reading, let $V$ be $2n$-dimensional with a non-degenerate alternating form, written as $U\oplus U^*$ with $U$ Lagrangian. A Lagrangian transverse to $U^*$ is the graph of a map $U\to U^*$, and the graph is isotropic if and only if the map is alternating. In that chart the intersection of the graph with $U$ is the kernel of the map. The uniform measure on transverse Lagrangians is the uniform measure on alternating matrices. The expectation of the size of the intersection is therefore the quantity just computed, and it tends to $p+1$.
\end{proof}

\begin{examplex}[$n=2$]
An alternating $2\times2$ matrix is $\begin{pmatrix}0&a\\-a&0\end{pmatrix}$. It is zero for one value of $a$, with kernel of size $p^2$, and invertible for the other $p-1$ values. The expectation is $p^{-1}\cdot p^2+(p-1)p^{-1}=p+1-p^{-1}$, in agreement with the theorem.
\end{examplex}

The prediction is that a $p$-Selmer group has average size $p+1$. It is a model. Bhargava and Shankar proved that the average size of the $2$-Selmer group, for elliptic curves ordered by height, is $3$, which is the case $p=2$ of the same number, and deduced that the average rank is at most $7/6$. That theorem is quoted. Goldfeld's conjecture, that half the quadratic twists have rank $0$ and half have rank $1$, is a conjecture.

\section{The congruent-number curve}

Volume~13 proves that $E_1:y^2=x^3-x$ has rank $0$. The map in that proof is the Kummer map, and its image is the $2$-Selmer group of this one curve.

The $2$-torsion is $\{O,(0,0),(1,0),(-1,0)\}$. An affine point $(x,y)$ with $y\ne0$ is sent to the pair of square classes of $x$ and $x-1$ in $\Q^\times/(\Q^\times)^2$. The image lands in the classes generated by $-1$ and $2$. The $2$-Selmer group is the set of those classes which lie in the local image at every place. For this curve the global computation of Volume~13 shows that the image is exactly the image of the torsion:
\[
\delta(O)=(1,1),\quad \delta(0,0)=(-1,-1),\quad \delta(1,0)=(1,2),\quad \delta(-1,0)=(-1,-2).
\]
The $2$-adic check is the one given there: a class with the factor $2$ in the first coordinate forces the valuation of $x$ to be odd, and then $x-1$ and $x+1$ cannot have even valuation simultaneously. Thus $\mathrm{Sel}_2(E_1)$ has order $4$. Since $E_1(\Q)/2E_1(\Q)$ injects into it and already has order $4$ from the torsion, the rank is $0$ and this descent sees no element of the Tate--Shafarevich group.

\begin{examplex}
The curve $y^2=x^3-25x$ has the rational point $(-4,6)$. Volume~13 turns it into a right triangle of area $5$. The same descent, with the prime $5$ allowed among the square classes, gives a Selmer group larger than the torsion, and the rank is positive.
\end{examplex}

The Birch--Swinnerton-Dyer conjecture, identifying the rank with the order of vanishing of $L(E,s)$ at $s=1$, is stated in Volume~13. Modularity, which makes that $L$-function entire, is a theorem, also stated there and not proved in this series. The meeting is the one this chapter can write down: the structure is the Selmer group, the randomness is a count of isotropic subspaces, and the sentence that they are the same object is a conjecture.
